\section{Methods}\label{sec:methods}

This section details our comprehensive methodological framework for analyzing magic squares through geometric and topological lenses. Our approach combines classical geometric analysis with modern computational topology techniques to reveal the intrinsic structure of magic square spaces.

\subsection{Geometric Embedding Framework}\label{subsec:embedding}

The foundation of our analysis rests on a natural embedding of magic squares into three-dimensional Euclidean space. For an $n \times n$ magic square $M$, we construct a point cloud in $\mathbb{R}^3$ by mapping each entry $M_{ij}$ to the point $(i, j, M_{ij})$, where $(i,j)$ represents the spatial position within the grid and $M_{ij}$ provides the height dimension.

This embedding preserves both the spatial relationships inherent in the square's structure and the numerical relationships between entries. Unlike alternative representations such as treating squares as vectors in $\mathbb{R}^{16}$ or using spectral embeddings, our approach maintains the geometric intuition of magic squares as balanced physical objects while enabling sophisticated topological analysis.

The choice of this embedding is theoretically motivated by recent work on manifolds with density \citep{bunge2020}. Magic squares can be viewed as discrete probability distributions over the grid positions, where the entry values serve as weights. Our embedding naturally captures this perspective, with the third coordinate providing the density function. This connection to geometric measure theory justifies our use of the Wasserstein (Earth Mover's) distance as the natural metric for comparing magic squares, as it respects both the spatial arrangement and the weight distribution. Furthermore, this framework allows us to apply results from optimal transport theory to understand the geometry of magic square spaces.

\subsubsection{Point Classification and Convex Hull Analysis}

For each embedded magic square, we compute its convex hull and classify points according to their geometric role:

\begin{itemize}
    \item \textbf{Vertex points}: Points that form vertices of the convex hull, representing extremal values in the geometric structure
    \item \textbf{Face points}: Points lying on faces of the convex hull but not at vertices
    \item \textbf{Interior points}: Points contained within the convex hull
\end{itemize}

Our analysis reveals that magic squares exhibit consistent patterns in point distribution. Table~\ref{tab:geom-stats} summarizes the geometric properties across all 880 unique $4 \times 4$ magic squares.

\input{sections/geometric_stats}

The geometric measures we compute include:
\begin{enumerate}
    \item \textbf{Hull volume}: The volume of the convex hull, measuring the geometric extent of the square
    \item \textbf{Surface area}: The surface area of the convex hull
    \item \textbf{Angular uniformity}: A measure of how evenly distributed the points are around the centroid
    \item \textbf{Compactness}: The ratio of volume to surface area, indicating geometric efficiency
\end{enumerate}

\subsection{Earth Mover's Distance Analysis}

To quantify similarity between magic squares, we employ the Earth Mover's Distance (EMD), also known as the Wasserstein distance. For two magic squares $M_1$ and $M_2$, the EMD measures the minimum cost of transforming the point cloud of $M_1$ into the point cloud of $M_2$.

Formally, given two point clouds $P_1 = \{p_1^{(1)}, \ldots, p_{16}^{(1)}\}$ and $P_2 = \{p_1^{(2)}, \ldots, p_{16}^{(2)}\}$ with uniform weights, the EMD is:

\[
\text{EMD}(P_1, P_2) = \min_{\pi \in S_{16}} \sum_{i=1}^{16} \|p_i^{(1)} - p_{\pi(i)}^{(2)}\|_2
\]

where $S_{16}$ is the set of all permutations of 16 elements, and $\|\cdot\|_2$ denotes the Euclidean norm in $\mathbb{R}^3$.

Our EMD analysis reveals natural clustering within the space of magic squares. Table~\ref{tab:emd-nn} presents statistics for nearest-neighbor distances.

\input{sections/emd_stats}

\subsection{Manifold Structure Analysis}

We investigate the manifold structure of magic squares using several complementary approaches:

\subsubsection{Local Dimension Estimation}

We estimate the local intrinsic dimension using the correlation dimension method. For each magic square, we compute the correlation integral:

\[
C(r) = \frac{1}{N^2} \sum_{i,j=1}^N \Theta(r - \|p_i - p_j\|)
\]

where $\Theta$ is the Heaviside step function, $N = 16$ is the number of points, and $p_i$ are the points in the embedded square. The local dimension is estimated as:

\[
d = \lim_{r \to 0} \frac{\log C(r)}{\log r}
\]

\subsubsection{Connectivity and Clustering Analysis}

We construct a $k$-nearest neighbor graph for the magic squares in EMD space and analyze its connectivity properties. Our analysis reveals the manifold structure summarized in Table~\ref{tab:manifold}.

\input{sections/manifold_stats}

\subsection{Persistent Homology Analysis}

To capture multi-scale topological features, we apply persistent homology to the point clouds of magic squares. For each square, we construct a filtered simplicial complex using the Vietoris-Rips construction:

\[
\text{VR}_\epsilon(P) = \{\sigma \subseteq P \mid \text{diam}(\sigma) \leq \epsilon\}
\]

where $P$ is the point cloud and $\epsilon$ is the filtration parameter.

The resulting persistence diagrams capture topological features (connected components, loops, voids) that persist across multiple scales. We compute persistence diagrams for dimensions 0, 1, and 2, corresponding to connected components, 1-dimensional holes, and 2-dimensional voids.

\subsubsection{Statistical Analysis of Persistence Features}

We develop statistical measures to characterize the persistence diagrams:

\begin{itemize}
    \item \textbf{Total persistence}: Sum of lifespans of all features
    \item \textbf{Maximum persistence}: Lifespan of the most persistent feature
    \item \textbf{Feature count}: Number of significant features above a threshold
    \item \textbf{Persistence entropy}: Information-theoretic measure of feature distribution
\end{itemize}

\subsection{Dual Null Model Framework}

To validate our findings, we compare magic squares against two null models:

\begin{enumerate}
    \item \textbf{Random permutation model}: Random arrangements of integers 1-16
    \item \textbf{Structured null model}: Arrangements satisfying some but not all magic square constraints
\end{enumerate}

This dual comparison framework allows us to distinguish features that are specific to magic squares from those that arise from general structural constraints.

\subsection{Computational Implementation}

All analyses are implemented in Python using the following key libraries:
\begin{itemize}
    \item \textbf{NumPy/SciPy}: Numerical computations and geometric algorithms
    \item \textbf{scikit-learn}: Manifold learning and dimension estimation
    \item \textbf{GUDHI}: Persistent homology computations
    \item \textbf{POT}: Optimal transport and EMD calculations
    \item \textbf{NetworkX}: Graph analysis for connectivity studies
\end{itemize}

The complete analysis pipeline processes all 880 unique $4 \times 4$ magic squares in approximately 30-40 seconds on modern hardware, making it suitable for interactive exploration and parameter studies.

\subsection{Validation and Reproducibility}

All results are validated through multiple independent implementations and cross-validation with analytical results where available. The complete codebase, data, and analysis scripts are provided as supplementary material to ensure full reproducibility of our findings. 